People searching together often find good answers faster than people searching apart, but a group that learns
from its members too readily can settle on a worse answer than its members would have found alone. Multi-agent systems built from large language
models (LLMs) usually let every agent see what its peers have tried and how it scored, so that nothing
is discovered twice. We ask whether that sharing helps or hurts when the score is an honest judge of the
current round but a poor guide to the best one, because reaching the best setup means scoring worse for
a while. We use a reimplementation of the lemonade-stand game of \citet{ederer2013pay}: teams of three
GPT-5.6 agents search hidden worlds under conditions that differ only in when a shared board of teammates'
moves and results is visible, with and without the original wording and its recommended starting setup. Three
results hold in both framings. \textbf{Sharing roughly halves the number of distinct setups a team tries.}
\textbf{Sharing raises what the team earns over the game.} \textbf{Sharing never reliably raises the best
single round any teammate reaches, and tends to lower it.} Agents that can see each other copy whoever is
doing best, which pays now and settles the team at a local optimum; the later the board opens, the more the
team explores. The first two results reappear on further worlds, a second model and a longer game; whether
the best round stays lower depends on the horizon and on a shared start. Sharing speeds the refinement of
the setup the team starts from; it does not help the team find out whether a better one exists.
